\documentclass[11pt,a4paper]{article}
\usepackage[utf8]{inputenc}
\usepackage[margin=1in]{geometry}
\usepackage{amsmath}
\usepackage{amssymb}
\usepackage{graphicx}
\usepackage{hyperref}
\usepackage{enumitem}
\usepackage{tcolorbox}
\usepackage{fancyhdr}

% Header setup
\pagestyle{fancy}
\fancyhf{}
\rhead{\textit{Aether Proactive System - Use Cases}}
\lhead{\textit{Research Implementation}}
\rfoot{\thepage}

\title{\textbf{Proactive Information Retrieval System:} \\
\large Real-World Use Cases from Multi-Source Activity Monitoring}
\author{AetherArena Implementation}
\date{January 2026}

\begin{document}

\maketitle

\begin{abstract}
We present a production implementation of proactive information retrieval (PIR) based on \textit{Generating Queries from User Contexts via In-Context Learning for Proactive Retrieval}. Our system operates in two phases: Phase 1 continuously monitors browser history (threshold: 3 logs), filesystem changes (threshold: 1 log), and email activity (threshold: 1 log) to generate 0-3 queries using pattern analysis with ICL (previous 3 batches as evolutionary context); Phase 2 employs a ReAct-based scouting agent with 8 specialized tools to assess relevance and make binary \texttt{intervene}/\texttt{defer} decisions. The LLM can choose to generate NO queries if activity is routine. We demonstrate realistic use cases across single-source and cross-source scenarios using actual data structures from production implementation.
\end{abstract}

\section{System Architecture Overview}

\subsection{Phase 1: Background Query Generation}
\textbf{Operates independently} of the main application—runs 24/7 as system daemons monitoring user activity across three primary sources:

\begin{enumerate}[leftmargin=*]
    \item \textbf{Browser Daemon:} Scans Chromium-based browser history every 5 seconds, capturing URL, title, visit count, and profile name. Uses incremental timestamp tracking to prevent duplicate processing.
    
    \item \textbf{Email Daemon:} Monitors macOS Mail.app via AppleScript (no Full Disk Access required), extracting subject, sender, and 500-character body preview from recent emails.
    
    \item \textbf{Filesystem Daemon:} Watchdog-based real-time monitoring of configured directories, logging file created/modified/deleted/moved events with 500-character content preview extraction for text files (40+ extensions supported: .py, .js, .md, .txt, .sql, etc.).
\end{enumerate}

When any source accumulates \texttt{threshold} unprocessed logs (browser: 3, email: 1, filesystem: 1), it signals the \textbf{Query Generation Daemon}, which:
\begin{itemize}
    \item Fetches unprocessed logs from all active sources
    \item Builds context window including current batch + last 3 batches (ICL continuity)
    \item Generates 0-3 queries via LLM using pattern analysis (paper Figure 3 methodology)
    \item LLM can choose to generate NO queries if activity is routine/non-interesting
    \item Stores queries in SQLite with bidirectional linking to source documents
    \item Updates BM25 index for agent retrieval
\end{itemize}

\subsection{Phase 2: Proactive Agent Worker}
\textbf{In-application worker} (10-second heartbeat) monitors Phase 1 query database:

\begin{enumerate}[leftmargin=*]
    \item Fetches most recent unprocessed query (single-query processing, marks older queries as stale)
    \item Calls ReAct agent with query + source documents + activity context
    \item Agent uses tool budget (max 2 calls per tool, max 10 total, early stop after 2 high signals):
    \begin{itemize}
        \item \texttt{retriever}: Unified local search across 6 sources
        \begin{itemize}
            \item BM25 mode: Daemon PyTerrier indexes (browser, email, filesystem, query\_gen)
            \item Semantic mode: Aether-RAG with FAISS-HNSW (browser, email, filesystem)
            \item \texttt{query\_gen}: Search Phase 1 generated queries + context\_doc\_ids (links to source logs)
            \item \texttt{proactive\_cache}: Past agent runs (Supabase semantic search)
            \item \texttt{chat}: Chat history (Aether-RAG)
        \end{itemize}
        \item \texttt{web\_search}: Live web search (DuckDuckGo/SearxNG)
        \item \texttt{scrape\_url}: Extract content from specific URLs
        \item \texttt{academic\_search}: Search academic papers and publications
        \item \texttt{social\_search}: Search Reddit, YouTube, and social platforms
        \item \texttt{legal\_search}: Search legal documents and case law
        \item \texttt{uploads\_search}: Search uploaded files and documents
        \item \texttt{done}: Signal completion of research phase
    \end{itemize}
    \item Decision logic: \texttt{intervene} if (highSignalCount $\geq$ 2) OR (highSignalCount = 1 AND timely AND actionable); else \texttt{defer}
    \item Relevance score (0-1) calculated for monitoring/logging, not primary decision factor
    \item If \texttt{intervene}: Generate response via LLM using gathered context from tool results
    \item Stores run in Supabase with decision, embeddings, and tool calls
    \item Emits WebSocket notification if \texttt{intervene}
\end{enumerate}

\section{Use Case Taxonomy}

We demonstrate two categories: \textbf{Single-Source} (activity in one daemon) and \textbf{Cross-Source} (correlated activity across multiple daemons with overlapping semantic content).

\section{Use Cases}

\subsection{Use Case 1: Online Shopping Price Comparison (Browser-Only)}

\textbf{Scenario:} User shopping for a laptop, visiting multiple retailer websites.

\textbf{Phase 1 Activity:}
\begin{tcolorbox}[colback=gray!5,colframe=gray!80,title=Browser Logs (3 visits in 5 minutes)]
\begin{verbatim}
[
  {
    "url": "bestbuy.com/laptops/dell-xps-15",
    "title": "Dell XPS 15 - $1,899 | Best Buy",
    "visit_count": 2,
    "profile": "Default",
    "timestamp": "2026-01-31T14:22:45Z"
  },
  {
    "url": "amazon.com/dp/B09XYZ/dell-xps-15",
    "title": "Dell XPS 15 Laptop - Amazon.com",
    "visit_count": 1,
    "timestamp": "2026-01-31T14:24:30Z"
  },
  {
    "url": "dell.com/en-us/shop/laptops/xps-15",
    "title": "XPS 15 Laptop | Dell USA Official",
    "visit_count": 1,
    "timestamp": "2026-01-31T14:27:10Z"
  }
]
\end{verbatim}
\end{tcolorbox}

\textbf{Phase 1 Output:}
\begin{tcolorbox}[colback=blue!5,colframe=blue!80]
\texttt{"dell xps 15 laptop price comparison best deal"}
\end{tcolorbox}

\textbf{Phase 2 Agent Execution:}
\begin{enumerate}[leftmargin=*]
    \item \textbf{Tool Calls:}
    \begin{itemize}
        \item \texttt{retriever(sources=["browser"], query="dell xps 15", mode="bm25")} → 8 past visits
        \item \texttt{web\_search("Dell XPS 15 best price 2026")} → 12 docs (price aggregators, student discounts)
        \item Early stop triggered (2 high relevance signals)
    \end{itemize}
    \item \textbf{Decision:} \texttt{intervene} (highSignalCount=2: \texttt{repeatedBrowserActivity} + high web search relevance)
\end{enumerate}

\textbf{Recommendation:}
\begin{tcolorbox}[colback=green!5,colframe=green!80]
\textit{I noticed you're shopping for a Dell XPS 15. Here's what I found:}
\begin{itemize}
    \item \textbf{Best Buy:} \$1,899 (in stock)
    \item \textbf{Amazon:} \$1,749 (Prime eligible, arrives in 2 days)
    \item \textbf{Dell Direct:} \$1,799 (student discount available: 10\% off = \$1,619)
\end{itemize}
\textit{Tip:} If you're a student, Dell's education store offers the best price. Also check if your credit card has price protection.
\end{tcolorbox}

\textbf{Outcome:} User saves \$280 by learning about student discount, making informed purchase decision.

\subsection{Use Case 2: Travel Planning from Email Confirmation (Email + Browser Cross-Source)}

\textbf{Scenario:} User receives hotel confirmation email, starts researching destination.

\textbf{Phase 1 Activity - Email (t=0):}
\begin{tcolorbox}[colback=gray!5,colframe=gray!80,title=Email Log]
\begin{verbatim}
{
  "subject": "Your Tokyo Stay is Confirmed - March 15-20, 2026",
  "sender": "reservations@hyatt.com",
  "body_preview": "Thank you for booking! Hyatt Regency Tokyo, 
                   5 nights, Check-in: March 15, 2026. Located in 
                   Shinjuku district. Confirmation #HTK-928475...",
  "file_path": "Mail.app",
  "timestamp": "2026-01-31T09:15:00Z"
}
\end{verbatim}
\end{tcolorbox}

\textbf{Phase 1 Activity - Browser (t+10min):}
\begin{tcolorbox}[colback=gray!5,colframe=gray!80,title=Browser Logs]
\begin{verbatim}
[
  {
    "url": "japan-guide.com/tokyo/shinjuku",
    "title": "Shinjuku Guide - Tokyo Travel",
    "timestamp": "2026-01-31T09:25:30Z"
  },
  {
    "url": "tripadvisor.com/restaurants-shinjuku-tokyo",
    "title": "Best Restaurants in Shinjuku",
    "timestamp": "2026-01-31T09:27:15Z"
  },
  {
    "url": "google.com/flights?departure=SFO&arrival=NRT",
    "title": "Flights from San Francisco to Tokyo",
    "timestamp": "2026-01-31T09:30:00Z"
  }
]
\end{verbatim}
\end{tcolorbox}

\textbf{Phase 1 Output:}
\begin{tcolorbox}[colback=blue!5,colframe=blue!80]
\texttt{"tokyo travel shinjuku restaurants things to do march 2026"}
\end{tcolorbox}

\textbf{Phase 2 Agent Execution:}
\begin{enumerate}[leftmargin=*]
    \item \textbf{Signal Detection:}
    \begin{itemize}
        \item \texttt{multiple\_source\_overlap=True} (email "Tokyo, March 15" + browser "Tokyo travel")
        \item \texttt{temporal\_keywords=True} ("March 15, 2026" - upcoming trip)
    \end{itemize}
    \item \textbf{Tool Calls:}
    \begin{itemize}
        \item \texttt{retriever(sources=["email"], query="Tokyo stay", mode="bm25")} → 1 doc
        \item \texttt{retriever(sources=["browser"], query="Shinjuku Tokyo", mode="bm25")} → 7 docs
        \item \texttt{web\_search("Tokyo travel tips March 2026")} → 15 docs
    \end{itemize}
    \item \textbf{Decision:} \texttt{intervene} (highSignalCount=2: \texttt{multiple\_source\_overlap} + \texttt{temporal\_keywords})
\end{enumerate}

\textbf{Recommendation:}
\begin{tcolorbox}[colback=green!5,colframe=green!80]
\textit{I see you're planning your Tokyo trip (March 15-20). Here's a starter checklist:}
\begin{itemize}
    \item \textbf{Weather:} Early spring, 10-15°C, pack layers + light jacket
    \item \textbf{Transit:} Get Suica/Pasmo card at airport for trains/subways (hotel is 5 min from Shinjuku Station)
    \item \textbf{Cherry blossoms:} Peak season! Shinjuku Gyoen Park (15 min walk from hotel) will be in full bloom
    \item \textbf{Don't miss:} Omoide Yokocho (food alley), Tokyo Metropolitan Govt Building (free observation deck)
\end{itemize}
\textit{Reminder:} Book flights soon for better prices (your search shows SFO→NRT).
\end{tcolorbox}

\textbf{Outcome:} User gets personalized travel prep checklist based on confirmed hotel location and dates.

\subsection{Use Case 3: Recipe Gathering for Dinner Party (Email + Filesystem Cross-Source)}

\textbf{Scenario:} User receives dinner party invitation email, saves recipes to local folder.

\textbf{Phase 1 Activity - Email (t=0):}
\begin{tcolorbox}[colback=gray!5,colframe=gray!80]
\begin{verbatim}
{
  "subject": "Dinner Party Next Saturday!",
  "sender": "maria@friends.com",
  "body_preview": "Hey! Hosting a small dinner party next Sat 
                   (Feb 8) at 7pm. Would love if you could 
                   bring an appetizer! Theme is Mediterranean. 
                   Let me know if you can make it!",
  "timestamp": "2026-01-31T15:30:00Z"
}
\end{verbatim}
\end{tcolorbox}

\textbf{Phase 1 Activity - Filesystem (t+30min):}
\begin{tcolorbox}[colback=gray!5,colframe=gray!80]
\begin{verbatim}
[
  {
    "action": "created",
    "file_path": "/Users/jane/Documents/Recipes/hummus_recipe.txt",
    "content_preview": "Mediterranean Hummus
                         Ingredients: chickpeas, tahini, lemon, 
                         garlic, olive oil, cumin
                         Prep time: 10 minutes...",
    "timestamp": "2026-01-31T16:00:00Z"
  },
  {
    "action": "created",
    "file_path": "/Users/jane/Documents/Recipes/bruschetta.txt",
    "content_preview": "Classic Bruschetta
                         Tomatoes, basil, garlic, balsamic, 
                         baguette slices
                         Serves: 8-10 people...",
    "timestamp": "2026-01-31T16:05:00Z"
  }
]
\end{verbatim}
\end{tcolorbox}

\textbf{Phase 1 Output:}
\begin{tcolorbox}[colback=blue!5,colframe=blue!80]
\texttt{"mediterranean appetizers recipes dinner party hummus bruschetta"}
\end{tcolorbox}

\textbf{Phase 2 Agent Execution:}
\begin{enumerate}[leftmargin=*]
    \item \textbf{Signals:} \texttt{multiple\_source\_overlap=True} (email "Mediterranean appetizer" + filesystem "hummus bruschetta recipes")
    \item \textbf{Tool Calls:}
    \begin{itemize}
        \item \texttt{retriever(sources=["email"], query="dinner party", mode="bm25")} → 1 doc (Maria's invitation)
        \item \texttt{retriever(sources=["filesystem"], query="recipes Mediterranean", mode="bm25")} → 2 docs (hummus + bruschetta)
        \item \texttt{web\_search("Mediterranean appetizers easy prep timeline")} → 10 docs (grocery, timing guides)
    \end{itemize}
    \item \textbf{Decision:} \texttt{intervene} (highSignalCount=2: \texttt{multiple\_source\_overlap} + \texttt{emailWithPublicUrls})
\end{enumerate}

\textbf{Recommendation:}
\begin{tcolorbox}[colback=green!5,colframe=green!80]
\textit{I see you're preparing for Maria's dinner party (Feb 8). You've saved hummus \& bruschetta recipes:}
\begin{itemize}
    \item \textbf{Grocery list:} Chickpeas (2 cans), tahini, tomatoes (6-8), fresh basil, baguette (2 loaves), olive oil, garlic, lemon
    \item \textbf{Timeline:} Both recipes are quick (10-15 min prep). You can make hummus 2 days ahead; bruschetta should be assembled fresh day-of
    \item \textbf{Serving tip:} Hummus serves 8-10, bruschetta serves 8-10—perfect for small party
\end{itemize}
\textit{Don't forget to reply to Maria's email!}
\end{tcolorbox}

\textbf{Outcome:} User gets shopping list + prep timeline generated from saved recipes and email context.



\section{Conclusion}

This report demonstrates three realistic use cases for our proactive information retrieval system: single-source browser activity (online shopping price comparison), cross-source email+browser correlation (travel planning from confirmation email), and email+filesystem integration (dinner party recipe preparation). Each use case illustrates the complete system flow from Phase 1 pattern-based query generation (0-3 queries with ICL from previous 3 batches) to Phase 2 ReAct-based agent scouting with tool budget constraints (max 2 calls per tool, max 10 total, early stop after 2 high signals). The system uses signal-based decision logic (highSignalCount $\geq$ 2 triggers intervention) rather than threshold-based relevance scoring.

All examples use actual tool signatures and data structures from production implementation.

\vspace{0.5cm}
\noindent\rule{\textwidth}{0.4pt}
\subsection*{Paper Alignment Note: Two-Agent ICL Architecture}

\textbf{Paper's Approach:} Uses rank-biased overlap (RBO) to precompute ICL query-pair examples from MS MARCO (8.8M docs, 400K queries), selecting reformulations with intermediate overlap (RBO $\sim$ 0.5) before query generation.

\textbf{Our Implementation:} Provides ICL to both agents via unified data flow:

\textbf{Phase 1 (Intra-Generation ICL):} LLM receives last 3 batches (query + source\_logs) in prompt. When generating Query N, sees:
\begin{itemize}
    \item \texttt{Batch N-2}: "laptop recommendations \$2000" + 3 logs (Reddit, YouTube, Best Buy)
    \item \texttt{Batch N-1}: "dell xps 15 student discount" + 3 logs (Dell site, reviews, sale email)
    \item \texttt{Current}: 3 new logs (checkout, credit card, specs PDF)
\end{itemize}

Stores complete context in SQLite (11 docs): current logs (3) + previous queries N-1, N-2 with their contexts (4 each).

\textbf{Phase 2 (Cross-Generation ICL):} Receives Query N with all 11 evolutionary context documents. Immediately sees progression (broad research → model selection → purchase prep) without tool calls. Can avoid redundancy (knows N-2 covered comparisons, N-1 covered Dell specifics). Has \texttt{retriever(query\_gen)} tool for deeper history beyond last 3 batches.

\textbf{Advantages:} Bidirectional ICL (both agents), no static corpus (uses user history), two-level memory (built-in + BM25 index), redundancy avoidance, on-demand expansion, real-time updates.

\noindent\rule{\textwidth}{0.4pt}

\end{document}
